\documentclass[a4paper,ngerman]{article}
\usepackage{amsmath}
\usepackage{amssymb}
\usepackage{babel}
\usepackage{color}
\usepackage{fancyvrb}
\usepackage[margin=2.5cm]{geometry}
\usepackage{listings}
\usepackage{marvosym}
\usepackage{url}
\usepackage{hyperref}
\setlength{\parskip}{1ex plus 0.2ex minus 0.2ex}
\setlength{\parindent}{0ex}
\DefineVerbatimEnvironment{CommandLine}{Verbatim}{fontsize=\footnotesize,frame=lines}
\newcommand{\zB}{z.\,B.}
\newcommand{\dH}{d.\,h.}
\newcommand{\caret}{$^{\wedge}$}
\lstset{
  defaultdialect=[gnu]make,
  defaultdialect=[gnu]c++,
%  numbersep=6pt,
%  literate={~}{$\sim$}1{^}{\caret}1,
  basicstyle=\footnotesize,
  basewidth=0.5165em,
  numbers=left,
  extendedchars=true,
  breaklines=true,
  breakindent=0in,
  prebreak={\mbox{{\color{red}\Righttorque}}},
  postbreak={\mbox{{\color{red}\Lefttorque}}},
  breakautoindent=false,
  frame=topbottom,
  aboveskip=\bigskipamount,
  belowskip=\bigskipamount,
  keepspaces=true,
  columns=[c]fixed,
  showstringspaces=true}
\begin{document}
\newcommand{\qtt}[1]{"`\texttt{#1}"'}
\title{Syntaxpr"ufer f"ur While- und Loop-Programme}
\author{Lehrstuhl f"ur Theoretische Informatik, Universit"at W"urzburg}
\maketitle
\section{Installation}
Zum Betrieb des Syntaxpr"ufers ist Python erforderlich. Abh"angig vom
Betriebssystem kann Python "uber die jeweilige Paketverwaltung oder
manuell installiert werden. Die offizielle Python-Webseite lautet
\url{http://python.org/}. Neben der genauen Sprachdefinition und
Referenz zur Standardbibliothek ist dort unter anderem die jeweils
aktuelle Python-Version zu finden.

Der Syntaxpr"ufer ben"otigt mindestens Python 2.6. Es wird jedoch
empfohlen, mindestens Python 3.0 zu installieren, da sich die
sp"ateren Erweiterungen von While- und Loop-Programmen auf
Sprachfeatures dieser Python-Version beziehen k"onnen.\footnote{Die
  Unterschiede zwischen Python 2.x und Python 3.x sind zum Teil
  betr"achtlich.}
\section{Bedienung}
Der Syntaxpr"ufer wird auf der Kommandozeile benutzt. Dazu ist dem
Aufruf von Python das Verzeichnis, in das der Syntaxpr"ufer entpackt
wurde, als Parameter mitzugeben. Alle dann nachfolgenden Parameter
betreffen den Syntaxpr"ufer selbst, wobei die erlaubten Parameter im
Folgenden erl"autert werden. Ein m"oglicher Aufruf des Syntaxpr"ufers
zusammen mit der resultierenden Ausgabe w"are zum Beispiel:
\begin{CommandLine}
goll@thomson:~$ python syntax-checker -e fac.txt 7
[+] <WHILE>  fac.txt is a while program.

[-] <LOOP>   fac.txt is not a loop program because of the following:
             Calling function `fac' from within its own definition not allowed. (line 12, column 15)

               l.12:      res = mul(fac((n - 1)), n)
                                    ^

I will now try and evaluate the function defined by fac.txt.

---[ start of print output ]--------------------------------------------
1
2
6
24
120
720
5040
---[ end of print output ]----------------------------------------------

The result is:
  main(7) = 5040
\end{CommandLine}
%$
Hier weist die an den Syntaxpr"ufer "ubergebene Option \qtt{-e} den
Syntaxpr"ufer an, das in \qtt{fac.txt} abgelegte Programm nicht nur auf While- und
Loop-Konformit"at zu pr"ufen, sondern auch die durch das Programm
definierte, im Beispiel 1-stellige Funktion an der Stelle 7
auszuwerten und das Ergebnis auszugeben.

Ruft man den Syntaxpr"ufer mit der Option \qtt{-h} auf, so wird eine
kurze Auflistung der erwarteten Parameter sowie der definierten Optionen
ausgegeben. Die erlaubten Optionen werden im Folgenden erl"autert.
\subsection{Allgemeine Optionen}
\begin{description}
\item[\texttt{--version}] Gibt die aktuelle Programmversion des
  Syntaxpr"ufers aus. Bei der Meldung von Fehlern im Syntaxpr"ufer ist
  es zur Reproduktion und Behebung des Fehlers n"otig, die verwendete
  Version des Programms zu identifizieren. Dies kann mit dieser Option
  erreicht werden.
\item[\texttt{-h}, \texttt{--help}] Zeigt die Kurz"ubersicht zur
  Bedienung des Syntaxpr"ufers an.
\item[\texttt{-v}, \texttt{--verbose}] Erh"oht den Detailgrad der
  w"ahrend des Programmablaufs ausgegebenen Meldungen; die Option kann
  mehrmals angegeben werden, um weitere Meldungen anzuzeigen. Im
  Normalfall gibt der Syntaxpr"ufer nur das Ergebnis "`erf"ullt While-
  bzw. Loop-Spezifikation"' mit etwaigen Fehlerstellen aus. Durch das
  Aktivieren von \qtt{--verbose} werden weitere Status- und
  Fortschrittsmeldungen sowie interne Datenstrukturen ausgegeben.
\end{description}
\subsection{Optionen zum Verhalten}
\begin{description}
\item[\texttt{-e}, \texttt{--evaluate}] Weist den Syntaxpr"ufer an,
  die durch das While- bzw. Loop-Programm definierte Funktion
  auszuwerten. Die Argumente der Funktion werden dabei nach dem
  Dateinamen angegeben, mit einem Parameter je Funktionsargument. Um
  \zB{} die durch das Programm \qtt{program.txt} definierte, 3-stellige
  Funktion an der Stelle $(1, 5, 7)$ auszuwerten, sind die Parameter
  \qtt{-e program.txt 1 5 7} anzugeben.
\item[\texttt{-t FILE}, \texttt{--test=FILE}] Testet die durch das
  Programm definierte Funktion mit Hilfe einer Testdatei
  (\qtt{FILE}). Eine Testdatei enth"alt Zeilen von Eingabetupeln und
  erwartetem Ausgabewert. Auf diese Weise k"onnen schnell und einfach
  viele Testf"alle durchgepr"uft werden. Eine Zeile der Form \qtt{1 5 7
  123} besagt dabei, dass die definierte, 3-stellige Funktion an der
  Stelle $(1, 5, 7)$ den Wert $123$ haben soll.
\item[\texttt{-T}] Aktiviert die print-Ausgabe auch bei Tests. Im
  Normalfall werden bei der Ausf"uhrung von Tests mit der Option
  \qtt{-t} die print-Ausgaben des While- bzw. Loop-Programms
  deaktiviert, um die Ausgabe "ubersichtlich zu halten. Mit dieser
  Option kann die Ausgabe der print-Anweisungen im Programm wieder
  aktiviert werden.
\end{description}
\subsection{Hinweise}
Aufgrund fundamentaler Probleme, aber auch bedingt durch Limitierungen
der Sprache Python, kann der Syntaxpr"ufer bei Auswertung von
Programmen nicht mit Sicherheit auf Nichtdefiniertheit pr"ufen. Im
Allgemeinen gilt: wenn der Syntaxpr"ufer bei Auswertung mit \qtt{-e}
oder Tests mit \qtt{-t} einen R"uckgabewert des Programms meldet, so ist
dieser korrekt. Es kann jedoch im Allgemeinen nicht von der Ausgabe
"`wahrscheinlich nicht definiert"' auf echte Nichtdefiniertheit
geschlossen werden,\footnote{Beispielsweise wird der
  Python-Interpreter ein tief rekursives Programm abbrechen, auch wenn
  es nach einigen weiteren Rekursionen einen definierten Wert annehmen
  w"urde. Die Begrenzung liegt hier darin, dass Python nur eine
  endliche Anzahl von geschachtelten Funktionsaufrufen erlaubt, um
  Endlosschleifen abfangen und behandeln zu k"onnen.} ebenso wenig
kann von einer scheinbar unendlichen Ausf"uhrung (Endlosschleife) auf
echte Nichtdefiniertheit geschlossen werden.\footnote{Dies ist die
  Essenz des sogenannten Halteproblems: es ist f"ur einen Algorithmus
  im Allgemeinen nicht m"oglich zu entscheiden, ob ein gegebener
  Algorithmus auf einer gegebenen Eingabe zu einem Ergebnis kommt,
  also anh"alt.}
\section{Beispiel}
Im Folgenden wird beschrieben, wie ein Programm erstellt und getestet
werden kann, das das Quadrat der Eingabe berechnet. Genaugenommen soll
das Programm die folgende 1-stellige Funktion berechnen:
\begin{equation*}
  \textrm{square}(n) = \begin{cases}
    n^2 & \text{falls $n > 0$,} \\
    0 & \text{sonst.}
  \end{cases}
\end{equation*}
Die Beschr"ankung auf positive Zahlen wird dabei die Implementierung
ein wenig vereinfachen.
\paragraph{Python installieren}
Wenn das verwendete Betriebssystem eine Paketverwaltung bereitstellt
(dies ist der Fall bei den meisten Linux-Varianten), so bietet es sich
an, damit Python 3.1 oder neuer zu installieren.\footnote{In Debian
  und Ubuntu ist dies beispielsweise das Paket \qtt{python3} oder
  alternativ \qtt{python3.1}. Pakete ohne expliziten Versionszusatz
  beziehen sich derzeit meist noch auf das "altere Python 2.x.} Ist
dies nicht m"oglich, so kann Python unter
\url{http://python.org/download/} f"ur das jeweilige Betriebssystem
heruntergeladen werden. Im Folgenden werden wir annehmen, dass Python
im Ausf"uhrungspfad installiert wurde, \dH, dass Python unter der
Kommandozeile direkt mit dem Befehl \qtt{python} aufgerufen werden kann.
\paragraph{Syntaxpr"ufer entpacken}
Der Syntaxpr"ufer wird in einem zip-Archiv namens
\qtt{syntax-checker.zip} angeboten.\footnote{Alternativ auch als
\texttt{gzip}-komprimierter \emph{tarball} unter dem Namen
\qtt{syntax-checker.tar.gz}.} Der Archivinhalt ist in ein Verzeichnis
mit dem Namen \qtt{syntax-checker} zu entpacken. Es bietet sich an,
dieses Verzeichnis an einem Ort zu platzieren, den man leicht auf der
Kommandozeile erreichen kann. Im Folgenden werden wir annehmen, dass
zur Arbeit mit dem Syntaxpr"ufer ein (anfangs leeres) Verzeichnis
namens \qtt{theoinf} angelegt wurde, in das wir im Unterverzeichnis
\qtt{syntax-checker} den Syntaxpr"ufer entpackt haben. In \qtt{theoinf}
werden wir auch unsere Programmdateien ablegen.
\paragraph{Programm schreiben}
Es ist an der Zeit, das eigentliche Programm zu schreiben, welches die
oben beschriebene Funktion \emph{square} berechnen wird. Dies kann in
einem beliebigen Texteditor geschehen. Wir werden das Programm in
einer Datei \qtt{tutorial.py} im Verzeichnis \qtt{theoinf} ablegen. Das
folgende Programm erf"ullt die gestellten Anforderungen:
\lstinputlisting[language=python,title=tutorial.py]{tutorial_1.py}
Da in While- und Loop-Programmen keine Multiplikation definiert ist,
mussten wir unsere eigene Implementierung programmieren. Es ist
festzuhalten, dass die Funktion \qtt{mul} nur f"ur Aufrufe mit $a \geq
0$ das korrekte Ergebnis der Multiplikation liefert, in allen anderen
F"allen wird 0 zur"uckgegeben. Dies trifft sich jedoch mit der
Definition von \emph{square}, weshalb wir in \qtt{square} ohne weitere
Fallunterscheidung $\mathrm{mul}(n, n)$ zur"uckgeben k"onnen.
\paragraph{Syntaxpr"ufer aufrufen}
In einer Konsole (auch: auf einem Terminal, auf der Kommandozeile)
k"onnen wir nun den Syntaxpr"ufer auf unser Programm \qtt{tutorial.py}
anwenden. Der folgende Aufruf liefert uns das erste Ergebnis:
\begin{CommandLine}
goll@thomson:~/theoinf$ python syntax-checker tutorial.py
[+] <WHILE>  tutorial.py is a while program.

[-] <LOOP>   tutorial.py is not a loop program because of the following:
             Function `mul' lacks initialization of variables. (line 1, column 5)

               l.1:  def mul(a, b):
                         ^
\end{CommandLine}
%$
Der Ausgabe zufolge handelt es sich also bei unserem Programm um ein
While-Programm, \dH, wir haben keine verbotenen (aber im "`normalen"'
Python definierten) Sprachkonstrukte verwendet. Allerdings gibt der
Syntaxpr"ufer auch aus, dass wir kein Loop-Programm vorliegen haben:
es fehle die Variableninitialisierung in der Funktion \qtt{mul} (und
auch in der Funktion \qtt{square}).
\paragraph{Nachbesserung}
Wir wollen aus unserem While-Programm durch geringf"uhige Anpassung ein
Loop-Programm machen. Dies hat den Vorteil, dass das Programm
garantiert terminiert, \dH, auf jede Eingabe wird unser Programm dann
eine definierte und eindeutige Ausgabe liefern. Dies ist zwar jetzt
bereits der Fall, kann aber dem Programm rein syntaktisch noch nicht
angesehen werden. Wir f"ugen die geforderten
Variableninitialisierungen hinzu und korrigieren \qtt{tutorial.py} somit
wie folgt:
\lstinputlisting[language=python,title=tutorial.py]{tutorial_2.py}
Wenn wir nun erneut den Syntaxpr"ufer ausf"uhren, erhalten wir die
gew"unschte Ausgabe:
\begin{CommandLine}
goll@thomson:~/theoinf$ python syntax-checker tutorial.py
[+] <LOOP>   tutorial.py is a loop program (and thus while program).
\end{CommandLine}
%$
\paragraph{Programm testen}
Bisher haben wir vom Syntaxpr"ufer nur die rein syntaktische Aufgabe
abverlangt, eine Eingabe auf While- bzw. Loop-Konformit"at zu
pr"ufen. Mit den Parametern \qtt{-e} und \qtt{-t} k"onnen wir die
definierte Funktion \emph{square} nun aber auch ganz konkret auf
bestimmte Eingaben auswerten. Der folgende Aufruf berechnet \zB{} das
Quadrat von 17:
\begin{CommandLine}
goll@thomson:~/theoinf$ python syntax-checker -e tutorial.py 17
[+] <LOOP>   tutorial.py is a loop program (and thus while program).

I will now try and evaluate the function defined by tutorial.py.

---[ start of print output ]--------------------------------------------
---[ end of print output ]----------------------------------------------

The result is:
  square(17) = 289
\end{CommandLine}
%$
Das Programm liefert das erwartete Ergebnis $289$ ($=17^2$). Bei der
Auswertung werden zwischen den Zeilen \qtt{start of print output} und
\qtt{end of print output} die Argumente von vorhandenen print-Anweisungen
ausgegeben. Dieses Hilfskonstrukt, welches die Berechenbarkeit nicht
beeinflusst, erlaubt eine elementare Fehlersuche.

Dar"uber hinaus k"onnen wir von einer weiteren F"ahigkeit des
Syntaxpr"ufers Gebrauch machen: wenn wir
beispielsweise eine Datei \qtt{test.txt} mit folgendem Inhalt anlegen,
k"onnen wir Tests auch automatisieren:
\lstinputlisting[title=test.txt]{test.txt}
Diese Datei enth"alt Eingabe- und erwartete Ausgabewerte der zu
testenden Funktion \emph{square}, und zu Demonstrationszwecken einen
offenbar falschen Test in der Zeile \qtt{999 0}. Ein Aufruf mit dem
Parameter \qtt{-t} testet nun alle hinterlegten Werte und gibt zu jedem
Test aus, ob dieser erfolgreich war:
\begin{CommandLine}
goll@thomson:~/theoinf$ python syntax-checker -t test.txt tutorial.py
[+] <LOOP>   tutorial.py is a loop program (and thus while program).

I will now try and test tutorial.py according to test.txt.

The test results are:

  [OK]  square(1) = 1
  [OK]  square(2) = 4
  [OK]  square(3) = 9
  [OK]  square(4) = 16
  [OK]  square(123) = 15129
  [  ]  square(999) = 998001, but expected 0
  [OK]  square(0) = 0
  [OK]  square(-1) = 0
  [OK]  square(-123) = 0

  Out of 9 tests, [OK] 8 passed and [  ] 1 failed.
\end{CommandLine}
%$
Die Ausgabe gibt Aufschluss dar"uber, welche Tests fehlgeschlagen
sind: in unserem Falle ist dies nur unser bewusst ung"ultiger
Test. Davon abgesehen funktioniert unser Programm korrekt und
berechnet (zumindest f"ur diese Beispielwerte) tats"achlich die
Quadratzahlen f"ur positive Eingaben und gibt ansonsten den Wert 0
zur"uck.
\subsection{Im CIP-Pool}
Im Folgenden wird kurz umrissen, wie die obige
Schritt-f"ur-Schritt-Anleitung im CIP-Pool der Mathematik und
Informatik in W"urzburg umgesetzt werden kann. Dabei sind die
genannten Befehle auf der Kommandozeile bzw. in einem Terminal (xterm
o.\,"a.) zu starten.
\paragraph{Python installieren}
Dies ist nicht n"otig, da Python auf den Rechnern des CIP-Pools
bereits vorinstalliert ist. Der Befehl \qtt{python} startet derzeit
Python 2.6.6, der Befehl \qtt{python3} startet Python 3.1.2. Je nach
Komplexit"at der Programme, die wir schreiben, und welche
Sprachfeatures sie verwenden, ist Python 3.x erforderlich. Wir werden
daher im Folgenden grunds"atzlich \qtt{python3} zum Aufruf verwenden.
\paragraph{Verzeichnis anlegen}
F"ur die folgenden Befehle werden wir uns ein eigenes Verzeichnis zur
Vorlesung im Heimatverzeichnis anlegen. Dort werden wir den
Syntaxpr"ufer sowie unsere eigenen Programme ablegen. Das Verzeichnis
m"ussen wir selbstverst"andlich nur einmal erstellen. Dies erreichen
wir wie folgt:
\begin{CommandLine}
s123456@aurich:~$ mkdir theoinf
\end{CommandLine}
%$
\paragraph{In Verzeichnis wechseln}
Damit die folgenden Anweisungen tats"achlich im neuen Verzeichnis
ausgef"uhrt werden, wechseln wir in dieses:
\begin{CommandLine}
s123456@aurich:~$ cd theoinf
\end{CommandLine}
%$
\paragraph{Syntaxpr"ufer herunterladen}
Da man sich zum Herunterladen des Syntaxpr"ufers bei WueCampus
anmelden muss, ist dies schneller in einem graphischen Browser
erledigt. Beispielsweise mit Iceweasel laden wir also den
Syntaxpr"ufer herunter (im tar.gz-Format) und speichern die Datei im
neuen Verzeichnis \qtt{theoinf} im Heimatverzeichnis.
\paragraph{Syntaxpr"ufer entpacken}
Nun entpacken wir das Archiv, das den Syntaxpr"ufer enth"alt:
\begin{CommandLine}
s123456@aurich:~/theoinf$ tar xzf syntax-checker.tar.gz
\end{CommandLine}
%$
Dadurch sollte ein neuer Ordner namens \qtt{syntax-checker} unter
\qtt{theoinf} hinzukommen.
\paragraph{Programm schreiben}
Unsere While- und Loop-Programme schreiben wir in einem Texteditor
unserer Wahl und platzieren sie im \qtt{theoinf}-Verzeichnis. Wir nehmen
an, wir h"atten das Programm \qtt{tutorial.py} auf diese Weise
reproduziert und in diesem Verzeichnis abgelegt.
\paragraph{Syntaxpr"ufer aufrufen}
Den Syntaxpr"ufer rufen wir aus dem Verzeichnis \qtt{theoinf} heraus wie
folgt auf und wenden ihn auf die dort hinterlegte Datei
\qtt{tutorial.py} an:
\begin{CommandLine}
s123456@aurich:~/theoinf$ python3 syntax-checker tutorial.py
[+] <WHILE>  tutorial.py is a while program.

[-] <LOOP>   tutorial.py is not a loop program because of the following:
             Function `mul' lacks initialization of variables. (line 1, column 5)

               l.1:  def mul(a, b):
                         ^
\end{CommandLine}
%$
Die anderen, weiter oben beschriebenen Syntaxpr"ufer-Befehle lassen
sich auf die gleiche Weise ausf"uhren. Parallel dazu kann nat"urlich
im Texteditor die Datei \qtt{tutorial.py} weiter bearbeitet und im
offenen Terminal immer wieder getestet werden.
\end{document}
